% Source: 40_Chapter_30_The_Rule_of_the_Saints_Puritanism.txt
\bookchapter[The Rule of the Saints]{CHAPTER 30}{THE RULE OF THE SAINTS}{Puritanism}

\leadin{I}{N 1630 MORE THAN} four hundred immigrants gathered at Southampton, England, preparing to
sail to the New World. John Cotton, a distinguished minister who would later join them on the other
side of the sea, preached a farewell sermon. His text for the occasion summed up the spirit of the
great adventure. He took it from 2 Samuel 7:10:

\begin{blockquote}
Moreover I will appoint a place for my people Israel, and will plant them, that they may dwell in a
place of their own, and move no more; neither shall the children of wickedness afflict them any
more, as before-time (KJV).
\end{blockquote}

Cotton declared that like the ancient Israelites these immigrants were God’s chosen people, headed
for the land he had promised and prepared. In this new land they would be able to labor undisturbed
for the glory of God.

There you have a snapshot of Puritanism: the Bible is the guidebook providing concrete instruction
for the believer’s conversion and the construction of a society designed to reinforce a biblical
view of the state; the church, often pictured analogously to ancient Israel, is also rightly
pictured in the Scriptures and should function as the agent for individual and societal reform; the
society and nation should also be ordered to bolster the purpose of God being played out in history.
While it endured as a distinct historic movement between 1560 and 1660, Puritanism provided for
Christians of every generation a model of the Christian faith as a decisive commitment to Jesus
Christ and how that life of the soul expresses itself in the public arena in a nation governed by
the truths of the Bible.

\section{PURITANISM: NEW LIFE AND NEW WORLD}

In more recent centuries Christians have renewed from time to time the Puritan experience of the
soul, a type of piety arising from the purely gracious work of God in the heart of the believer. But
no later Christian movement has had the vision or the opportunity to apply the law of God to the
life of a nation.

In modern times, marked by zeal for individual rights and sexual freedom, puritan has come to refer
to a person who is self-righteous, hyper-religious, or sexually repressed: one who does his best to
keep people from having fun. This view of the puritan as a moss-backed moralist captured popular
thinking in the backlash of Victorian stuffiness. Early in the twentieth century American journalist
H. L. Mencken summed up the popular image in his quip that Puritanism was “the haunting fear that
someone, somewhere may be happy.”

But is that fair? What was Puritanism originally? Whatever it was, it was not straitlaced. It stood
for change and a new day in England. The first Puritans had little confidence in traditional
religion. Their plans for a new England arose from a deep conviction that spiritual conversion was
crucial to Christianity. This rebirth separated the Puritan from the mass of mankind and endowed him
with the privileges and the duties of the elect of God. The church may prepare a man for this
experience, and, after it, the church may guide him, but the heart of the experience, the reception
of the grace of God, is beyond the church’s control.

In its quest to reshape England, the Puritan movement passed through three rather clearly marked
periods: First, under Queen Elizabeth (1558–1603) it tried to purify the Church of England along the
lines of Calvin’s Geneva. Second, under James I and Charles I (1603–1642) it resisted the claims of
the monarchy and suffered under royal pressures designed to force conformity to a high-church style
Christianity. Third, during England’s Civil War and Oliver Cromwell’s rule (1642–1660), Puritans had
a chance to shape the national church in England but failed because of their own internal
dissensions.

Such an outline suggests, as it should, that Puritanism had a public side as well as a personal
side. It began with an individual’s experience of the redeeming grace of God but moved on to stress
the elect’s mission in the world, the shaping of society according to biblical principles.

In its emphasis upon the interior life of the saint, Puritanism was a taproot of later evangelical
Christianity with its born-again message. In its stress upon a disciplined “nation under God” and
his laws, it contributed significantly to the national character of the American people.

Puritanism, England’s second reformation, first appeared during the reign of Elizabeth. It had a new
style of preaching: a message aimed not at the head but at the heart. This “spiritual brotherhood,”
as William Haller the Puritan authority once called it, included names like Greenham, Rogers,
Chaderton, and Dod. Many of these first Puritans had been Protestant exiles of the reign of Bloody
Mary (1553–1557). Hounded out of their homeland by their Catholic queen, these Protestant
sympathizers had gone to Geneva and mobilized as a vanguard for a fresh Calvinist assault on
England.

The military image is appropriate, for in the late 1550s Geneva had become a center for
international subversion—a Protestant Moscow, if you please—sending out cadres of students enflamed
with a passion to overthrow Catholicism in their homelands. In 1560 John Knox was successful in
Scotland, and the English Reformers hoped for as much in their own country.

Queen Mary’s death and the crowning of Elizabeth appeared to be an invitation to the exiles to
return to England. They took up their work, still convinced that the Bible gave specific guidance
for the ordering of personal life, the structuring of the church, and the regulation of society.
They discovered that the most explosive political issue was the authority to choose ministers for
the churches. The Puritans wanted the selection of their ministers to rest with the people; the
queen insisted that the appointment of bishops was a responsibility of rulers.

As long as Elizabeth lived she allowed these Puritan dogs to bark so long as they did not bite. The
queen’s tolerant policy permitted them to complain about prayers from a book and special clothing
for ministers and the sign of the cross during baptisms—just about any practice that looked like
Roman Catholicism—but she would not hear of their control of the Church. She allowed them to lecture
for hours on the importance of preaching and the biblical office of elder, but she made sure the
Church of England remained firmly in the hands of bishops, and therefore in the control of the
crown.

\section{A PEOPLE OF THE BOOK}

The Puritan’s zeal to purify the Church of England was fired by eager reading of popular versions of
the Bible. Chief among these was the Geneva Bible. The name came from the fact that it was the work
of several of those English exiles in Geneva during the reign of the Catholic Mary. Miles Coverdale,
who had translated the first complete Bible into English in 1535, was in this group. Once printed
during the early years of Elizabeth’s reign, the Geneva Bible’s numbered verses, lucid prose,
improved scholarship, extensive prologues, and marginal notes gave it wide popularity. Until its
eventual replacement by the King James Version (1611) it was the most widely distributed Bible in
English and was the one the Puritans carried to America.

The Puritans, however, had more than their Geneva Bibles. They had a sense of destiny, a vision of
God’s purpose among men and nations. The idea that the Almighty moves in the affairs of men is
traceable, to be sure, to the Bible. But the concept that the pilgrimage of God’s people spanned the
centuries and led at last to England was from another book. John Foxe, in his highly popular Book of
Martyrs, planted this idea in the minds of Englishmen.

Like the Geneva Bible, The Book of Martyrs was a product of the English exile during the reign of
Bloody Mary. Foxe marshaled account after account of the suffering of faithful Protestants who dared
to die, if need be, for the triumph of God’s kingdom. According to Foxe, this trail of martyrs led
to the shores of England and to the reign of Mary. The conclusion seemed clear: God had a special
place for the English people in his worldwide plan of redemption.

The influence of The Book of Martyrs proved enormous. Generation after generation of Englishmen saw
history, and read their Bibles, through Foxe’s eyes. Aside from the Bible, his book probably did
more to shape the mind of Englishmen than any other single volume. Less than half a century after
Foxe’s death, Puritans carried his work and his philosophy of history, along with their English
Bibles, across the Atlantic into the American wilderness.

Thanks to their study of the Bible and their reading of The Book of Martyrs, Puritans came to think
of themselves as God’s new Israel. Most Christians in the seventeenth century believed that the
Bible was God’s inspired Word and revealed will. The key to the Puritan view of the Bible and of
themselves lies in their understanding of that fundamental biblical concept called covenant.

The Puritans, like the ancient Hebrews, believed that spiritual contracts exist between God and men.
The most fundamental was the covenant of grace. This was the spiritual contract by which true
Christians were bound to God. While they believed that God sovereignly elected men to salvation,
they also held that anyone by personal faith in Jesus Christ could be added to the covenanted
company. By grace believers became God’s people and he their God.

This “near bond” obligated the saint to walk in all of God’s ways made known in his Word. On the
other hand, it opened the Scriptures as a deeply felt source of spiritual and emotional strength, a
fountain of that rugged determination for which the Puritan became famous. To live within the
covenant of grace was to live within the light of the Word and according to the plan of Almighty
God.

\section{CONFLICTS WITH THE KING}

In 1603 the long reign of Queen Elizabeth came to a close when she died without an heir. James VI of
Scotland, the son of Mary Queen of Scots, became James I of England, uniting for the first time the
two kingdoms. Any Puritan hopes for James bringing Presbyterianism from Scotland to England were
dashed early. He welcomed a chance to deal with bishops. Rule in Scotland had been a constant
struggle with Presbyterian ministers. A Scottish presbytery, he said, “agrees with monarchy as well
as God and the devil.”

At the Hampton Court Conference in 1604 some leading Puritans had a chance to present to the king
their ideas for change in the Church of England. But James, who had an inflated view of his own
intelligence, dismissed most of their opinions rather rudely. On only one point did he consent to
the demands of the Puritans. He was willing to have a new translation of the Scriptures made. From
this decision came what we call the King James Version of the Bible.

On nothing else of significance would the king yield to the Puritans. The ceremonies, the Prayer
Book, and the bishops of the Church of England were going to stay. And if the Puritans did not like
it, they would still submit or, said the king, “I will drive them out of this land—or else worse.”
With that threat the conference broke up.

James’s initial contact with Parliament was equally tactless. In his first address to the body he
defended the divine right of kings. “The state of monarchy,” he said, “is the supremest thing upon
earth, for kings are not only God’s lieutenants upon earth and set upon God’s throne, but even by
God himself they are called gods.”

Sweeping aside English institutions and ignoring the mood of the English people, James made it plain
that he meant to be an absolute monarch. In 1611 he dissolved Parliament and for the next ten years
he ruled England without it. Thus the leaders of the Puritans and the advocates of parliamentary
authority in England tended to merge in their resistance to royal power.

Some in the Puritan movement, however, grew impatient for change in the Church. Shortly after the
Hampton Court Conference little groups of believers began to meet for worship as they felt the Bible
taught them, not according to bishops and prayer books. They were determined to obey God even if
their nation’s leaders were not. We call this movement Separatism because the groups were intent
upon leaving (to build a separate church) the Church of England.

One of these groups was in the village of Scrooby in northern England. Another was not far away at
Gainsborough. By 1608 both congregations had moved to Holland for safety and freedom of worship. The
Scrooby group settled at Leyden; the Gainsborough group found a home in Amsterdam.

After ten years in Holland the Scrooby congregation, led by Pastor John Robinson, realized that
their children were growing up out of touch with their homeland, forgetting even their native
tongue. But to return to England meant a return to the evil ways of English society. They had heard
of the English colony in Virginia planted in 1607. Perhaps the new continent of America was their
answer. It was a daring thought, full of unknown terrors, but a small number was willing to try.
These returned to England, where they were joined by another group of like-minded pilgrims. In
September 1620 a company of about a hundred souls set sail from Plymouth in a ship called the
Mayflower. In November the crew sighted the rocky shore of New England. The Pilgrim Fathers stepped
ashore at what we call Plymouth, Massachusetts.

Meanwhile, the second group sought safety from persecution in Holland. These separatists had been
given aid by local Mennonites. We know little about their association, but we know they exchanged
confessions to compare doctrines. The pastor of the flock, a Cambridge graduate, John Smyth, studied
his Greek New Testament and discovered that the practice of baptizing babies never appeared in its
pages. If babies were not included in the covenant of grace, only mature believers in Jesus Christ,
then shouldn’t churches be constituted by confession of faith rather than ties of covenants? Upon
their profession of faith they constituted the first Baptist church. Smyth baptized himself and then
the other forty members of his congregation by pouring. The year was 1609. The question of Baptist
origins and identity are closely tied. Are Baptists to be understood as English Separatists or
English Separatists who mingled with Anabaptists?

\section{ATTACK UPON THE LORD’S ANOINTED}

Separation from the Church of England, however, was far too drastic for the vast majority of
Puritans. They remained hopeful for some alternative to schism. All grounds for hope in earthly
rulers appeared to vanish in 1625, however, when Charles I followed his father James to the English
throne. James had his theories about the divine right of kings; Charles was determined to put them
into practice. No law and no parliament would restrain him. To make matters worse in the eyes of the
Puritans, his wife, Henrietta Maria, was a French Roman Catholic princess.

To enforce his will upon the Puritans, Charles found a ready servant in Archbishop William Laud, who
believed that God had ordained bishops to govern his church. Led by the archbishop, an episcopal
party arose to resist the Puritans. With the king’s support they reintroduced stained glass windows,
crosses—even crucifixes. They elevated the Communion table and called it an altar and they insisted
that worship be conducted according to the Prayer Book and no other.

Laud’s high-minded and high-handed policy drove some Puritans toward Separatism and others across
the Atlantic to America. Within ten years after Laud became archbishop, twenty towns and churches
had sprung up in Massachusetts Bay: in all sixteen thousand people including the four hundred who
heard John Cotton’s farewell at Southampton.

The straw that broke the back of restraint within the two kingdoms came when Charles tried to force
his high church brand of Anglican religion upon the Presbyterian Scots. He insisted that they
conduct their worship services from the Book of Common Prayer. John Milton called it “the skeleton
of a Mass-book.” That is the way the Scots took it. They rose in opposition to the move and joined
in a “National League of Covenant.” To defend their Church, they dared to take up arms against their
king.

To put an army in the field Charles was forced to convene Parliament, which he, like his father, had
tried to ignore for over a decade. Once Parliament assembled, conflicting loyalties led to a
Royalist Party and a Parliamentary Party. The Parliamentarians, clearly a majority, were agreed on
the broad Puritan principles but were divided over the form of the church. On the one hand were
Presbyterians; on the other were Independents (or Congregationalists). United in their hatred of
Archbishop Laud, the Parliamentarians succeeded in bringing him to trial and seeing him beheaded.

When Charles tried to punish the leaders of this opposition, civil war erupted. The Royalist members
of Parliament left London to join the forces defending the king. So Parliament was free at last to
institute the reform of the Church the Puritans had always wanted. It called to Westminster scores
of Puritan theologians and assigned to them the creation of a new form of worship and a new form of
church government for the Church of England.

The Westminster Assembly, meeting from 1643 to 1649, produced the Westminster Confession of Faith to
replace the Thirty-nine Articles, as well as a Larger and Shorter Catechism for use in the churches.
These writings alone made the assembly one of the most significant gatherings of Christian history.
Many orthodox Presbyterians and Congregationalists use these documents to this day.

In 1645 Parliament ordered the creation of presbyteries throughout the land and the election of
elders for the churches. The system, however, was never fully in force because Parliament never
ruled the entire country, and even when the king was defeated, the Army, not Parliament, wielded the
power.

\section{THE RULE OF THE SAINTS}

The war itself brought forward one of the towering figures of Christian history, a country gentleman
named Oliver Cromwell (1599–1658). As a colonel in the Parliamentary forces he proved to be a
military genius. His regiment, known as Ironsides, was never defeated, in part because Cromwell
instilled in his men a sense of discipline and Christian mission.

The rise of Cromwell’s star brought him in time to the leadership of the New Model Army, a force of
twenty-one thousand men, who considered their role in English history a call from God. In their
eyes, the war was a Puritan crusade against the enemies of righteousness. The just end sanctified
force as a means. That is why the Army prayed before battles and plunged into conflict singing
Christian hymns.

By the end of 1646 Cromwell’s Army had forced Charles to surrender. For the next two years the king
tried to play off his enemies against each other: the Scots, the Presbyterians (who controlled
Parliament), and the Independents (who dominated the Army). He succeeded in splitting Parliament and
making a secret alliance with the Scots. But fierce resentment against Charles broke out in the
Army, and in 1648 war erupted anew.

This time the Army defeated the allies of the king, and the Presbyterians were purged from the House
of Commons. This so-called Rump Parliament, a tool of the Army, created a high court of justice to
try the king. In January 1649 Charles was led to the scaffold in front of the royal palace of
Whitehall in London, and, before an assembled throng, he was executed.

It was a brash move, a sure sign of Puritanism’s ultimate fall from power, for it gave the royalists
their own martyr. Centuries of English royal tradition could not be erased, even by the saints of
God.

\figtitle{The Reformation in England}

Shortly after the king’s execution a portrait appeared, allegedly depicting his last hours. He was
kneeling at a table; on it rested the Bible. The royal crown lay on the floor. In the king’s right
hand was a crown of thorns. His eyes were lifted to a crown of glory above.

The widely distributed picture and its attendant sentiment unmasked Puritanism as no battlefield
could. It changed Charles from a king executed by rebels into the martyr persecuted by fanatics.
That is how many Englishmen chose to remember the Puritans.

For the moment, however, the Army of saints had their say. After abolishing the House of Lords, the
House of Commons proclaimed England a republic: the Commonwealth. But in 1653 the Army, still
distrusting Parliament, overthrew the Commonwealth and set up a form of government called the
Protectorate. Cromwell held the Office of Lord Protector, virtually a military dictator of England.

The Lord Protector tried to achieve a religious settlement for the nation by granting liberty to a
wide variety of Christian groups now growing on the religious landscape: Presbyterians,
Independents, Baptists, Quakers, Levellers, and others. Unfortunately he found the task impossible,
and the last three years of his life were filled with disappointment and trouble. When he died in
1658 the rule of the saints in old England died with him. Within two years the country welcomed the
return of the monarchy, and with the king the office of bishop. Far across the Atlantic, however,
tens of thousands of saints were happily settled in New England where they were devoutly erecting
the kingdom of God in the American wilderness.

\readinghead{Suggestions for Further Reading}
\begin{readinglist}
\entry *Bebbington, David. Baptists through the Centuries: A History of a Global People. Waco, TX:
  Baylor Press, 2010.
\entry *Brackney, William H. The Baptists. Westport, CT: Praeger, 1995.
\entry Haller, William. The Rise of Puritanism. New York: Harper \& Brothers, 1938.
\entry Morgan, Edmund S. Visible Saints: The History of a Puritan Idea. Ithaca, NY: Cornell
  University Press, 1963.
\entry Rutman, Darrett B. American Puritanism. Philadelphia: Lippincott, 1970.
\entry Simpson, Alan. Puritanism in Old and New England. Chicago: University of Chicago Press, 1955.
\entry Torbet, Robert G. A History of the Baptists. Philadelphia: Judson Press, 1963.
\entry Watkins. Owen C. The Puritan Experience. New York: Schocken, 1972.
\end{readinglist}
